% =====================================================================
% CAPITULO 4 - LA EVALUACION INTRA-DIA
% =====================================================================
% TESIS: dentro del dia el problema esta saturado y la evaluacion intra-dia no
% discrimina entre vistas. La causa que se propuso primero (un solo host
% atacante por dia -> memorizacion del host) la refuta el leave-one-attacker-out
% de la seccion 4.4. El mecanismo que queda es la herramienta, y lo confirma la
% evasion. La discriminacion real se traslada al capitulo 5.
%
% El marco de los cuatro regimenes se conserva como descripcion de donde vive
% la firma de cada ataque, no como asignacion de vista ganadora.
%
% CONVENIO DE TABLAS: flotante `table` con \caption arriba, \label despues del
% caption y booktabs. Citada desde el texto como "la tabla~\ref{...}".
%
% METRICA: macro-F1 out-of-fold + recall de ataque, con IC 95% cuando N es
% pequeno. No citar exactitud: la seccion 4.1 explica por que engana.
%
% ESTILO (31-ago): sin negritas en el texto corrido, sin incisos entre rayas
% salvo cuando la frase los pida, y sin florituras.
% =====================================================================

Este capítulo mide. Toma los seis días procesados del CSE-CIC-IDS2018, aplica a
todos el mismo protocolo y las mismas tres representaciones de los datos, y
compara. El resultado no es el que se esperaba al empezar, y es el hallazgo
central del trabajo: dentro de un mismo día el problema está saturado.
Cualquiera de las vistas detecta cualquiera de los ataques con un macro-F1 por
encima de $0.97$, de modo que la comparación no distingue entre ellas.

Eso no es un fracaso de la medición sino un resultado sobre el conjunto de
datos, y tiene una causa identificable que la sección~\ref{sec:por-que-satura}
documenta. Su consecuencia ordena el resto de la memoria: si la evaluación
dentro del día no discrimina, la pregunta de qué vista es mejor hay que
plantearla en otro sitio. Ese otro sitio, la robustez ante evasión, la
transferencia entre dominios y la calibración, es el
capítulo~\ref{ch:generalizacion}.

\section{Marco de comparación}
\label{sec:marco}

Todas las comparaciones de este capítulo y del siguiente siguen el mismo
protocolo, que se fija aquí para no repetirlo. Las decisiones que lo componen no
son neutrales: cada una corrige un sesgo concreto que se detectó
experimentalmente durante el trabajo, y sin ellas los resultados serían
sistemáticamente optimistas.

\subsection*{Selección por servicio, no por puerto}

El subconjunto de flujos sobre el que se evalúa cada día se define por el
servicio identificado por contenido, con las mismas firmas que emplean los
analizadores de Zeek: el \emph{banner} SSH, los métodos HTTP, el registro TLS o
los saludos SMTP y FTP. No se define por el número de puerto. El puerto depende
de la configuración de cada red y no es, por tanto, un atributo independiente
del contexto: un mismo servicio puede escucharse en puertos distintos y un mismo
puerto alojar servicios distintos. Usarlo como criterio ata los resultados a la
topología del laboratorio.

La corrección no alteró ninguna conclusión previa, porque en un entorno
controlado el contenido confirma al puerto casi siempre, pero sí cambia lo que
se puede afirmar. Y cambia algunos recuentos: en el día del SSH-Bruteforce el
tráfico benigno genuino son 433 flujos si se selecciona por el puerto 22 y 358
si se selecciona por contenido. Las cifras que se reportan a partir de aquí son
siempre las de la selección por contenido.

\subsection*{Un flujo es una conexión TCP}

La unidad de ejemplo es la conexión TCP, identificada por el \texttt{uid} que
asigna Zeek, y no el paquete. Un ataque se manifiesta a lo largo de una
conversación, no en un datagrama aislado, y agregar por conexión evita que un
mismo evento aparezca centenares de veces en el conjunto de entrenamiento y una
sola vez en el de prueba.

\subsection*{Balanceo 1:1, validación cruzada y repeticiones}

El tráfico benigno representa más del 95\% de los flujos en todos los días
analizados, así que la exactitud global es una métrica engañosa: un
clasificador que responda siempre «benigno» la maximiza. Cada comparación se
hace sobre un subconjunto balanceado 1:1 por submuestreo aleatorio de la clase
mayoritaria, con un tope de 5.000 ejemplos por clase. Con ese balanceo el azar
corresponde exactamente a $0.5$.

La evaluación es validación cruzada estratificada de cinco particiones con
predicciones \emph{out-of-fold}: cada flujo se predice una única vez, por el
modelo que no lo vio durante el entrenamiento.

El clasificador es el mismo en todas las comparaciones de este capítulo: un
perceptrón multicapa de dos capas ocultas, de 256 y 128 unidades, precedido de
un escalado estándar que se ajusta únicamente sobre la partición de
entrenamiento. La sección~\ref{sec:modelos} lo describe junto con el resto de
modelos del trabajo. Mantenerlo fijo es deliberado, porque lo que aquí se
compara son las representaciones y no las arquitecturas: si el modelo cambiase
entre vistas, la diferencia medida no sería atribuible a la información que cada
vista aporta. Las arquitecturas se comparan entre sí en el
capítulo~\ref{ch:generalizacion}, y allí lo que se mantiene fijo es la vista.

Hay un detalle del submuestreo que afecta a la
reproducibilidad. En el día del SSH el balanceo elige 358 flujos de ataque de
entre 94.207, y cuáles se elijan mueve el resultado en la tercera cifra decimal,
que es justo donde se separan las vistas. Fijar una semilla arbitraria y
reportar ese número sería presentar como exacto algo que no lo es. Por eso cada
comparación se repite diez veces con semillas distintas y se reporta la media.

Esas diez repeticiones no son diez
semillas de la validación cruzada. Cada repetición extrae un subconjunto
balanceado distinto y lo evalúa entero con validación cruzada de cinco
particiones: el procedimiento son diez ejecuciones completas de la
validación cruzada, es decir, cincuenta ajustes del modelo por vista y día. Lo
que varía entre repeticiones es qué flujos de la clase mayoritaria entran en la
comparación, que es la fuente de incertidumbre que importa cuando el balanceo
descarta la mayor parte de esa clase, como aquí, que reduce los 94.207 flujos de
ataque a los 358 benignos genuinos disponibles.

El algoritmo~\ref{alg:protocolo} recoge el procedimiento completo, porque
descrito en prosa queda repartido entre varios párrafos y su reproducción exacta
importa. Las dos fuentes de incertidumbre que el intervalo recoge corresponden a
los dos bucles: el exterior sobre submuestreos y el interior sobre remuestreos.

\begin{algorithm}[htbp]
\caption{Protocolo de comparación entre vistas}
\label{alg:protocolo}
\begin{algorithmic}[1]
\Require Flujos de ataque $A$ y benignos genuinos $B$ del servicio; vistas $V$;
         repeticiones $R = 10$; remuestreos $N = 2000$
\Ensure Macro-F1 medio e intervalo de confianza al 95\% por vista
\State $n \gets \min(|A|, |B|, 5000)$
\ForAll{vista $v \in V$}
  \State $P_v \gets \emptyset$;\ \ $Q_v \gets \emptyset$
  \For{$r = 1$ \textbf{hasta} $R$}
    \State $S \gets n$ flujos de $A$ y $n$ de $B$, al azar con semilla $r$
    \State $\hat{y} \gets$ predicción \emph{out-of-fold} de $S$ con 5-Fold estratificado
    \State $P_v \gets P_v \cup \{\text{macro-F1}(y_S, \hat{y})\}$
    \For{$b = 1$ \textbf{hasta} $N / R$}
      \State $I \gets$ muestra con reemplazo de los índices de $S$
      \State $Q_v \gets Q_v \cup \{\text{macro-F1}(y_S[I], \hat{y}[I])\}$
    \EndFor
  \EndFor
  \State \textbf{reportar} media de $P_v$ y percentiles 2.5 y 97.5 de $Q_v$
\EndFor
\end{algorithmic}
\end{algorithm}

\subsection*{Benigno genuino: el error que se cometió dos veces}
\label{sec:benigno-genuino}

La decisión más importante del protocolo es también la que más costó aprender.
El tráfico benigno de referencia debe proceder de terceros, nunca del propio
atacante fuera de la ventana etiquetada.

Parece obvio enunciado así, pero el error se cometió dos veces
(secciones~\ref{sec:benigno} y~\ref{sec:reg-volumetrico}). En el día del
SSH-Bruteforce, de los 2.022 flujos del puerto 22 etiquetados como benignos,
1.589 eran del propio atacante, conexiones suyas que caían fuera de los límites
de la ventana, y solo 433 procedían de terceros. Con ese conjunto se estaba
comparando ataque contra ataque, y el resultado era una medida sin sentido que
además parecía razonable: ninguna vista superaba el $0.78$ y se interpretó
como el techo del problema. En el día del DoS-Hulk el mismo error afectaba al
65.1\% del benigno HTTP.

La corrección no consistió en filtrar esos flujos al evaluar, sino en arreglar
la ventana y reetiquetarlos en el conjunto de datos, porque eran actividad del
ataque. Tiene una consecuencia metodológica que declarar: el escenario
contaminado ya no puede reconstruirse desde los datos actuales, de modo que las
cifras que en su día lo describieron son históricas y no comparables con las de
este capítulo.

Lo que generaliza es la regla. Hay que verificar que la ventana etiquetada cubre
toda la actividad continua del atacante sobre el conjunto completo, y no sobre
una captura suelta. Un artefacto de etiquetado no se detecta mirando las
métricas, porque las del escenario contaminado eran perfectamente plausibles,
sino auditando de dónde procede cada flujo de la clase negativa.

\subsection*{Métricas}

Se reportan macro-F1 y \emph{recall} de la clase de ataque. La primera resume el
equilibrio entre precisión y sensibilidad sin dejarse arrastrar por la clase
mayoritaria. La segunda responde a la pregunta operativa de cuántos ataques se
escapan, y se reporta siempre por separado porque un macro-F1 de $0.33$ puede
corresponder tanto a un clasificador mediocre como a uno que no detecta
absolutamente nada.

Ambas van acompañadas de un intervalo de confianza al 95\% obtenido por
\emph{bootstrap} de 2.000 remuestreos repartidos entre las diez repeticiones del
submuestreo, de modo que el intervalo recoge las dos fuentes de incertidumbre:
el tamaño finito del conjunto evaluado y qué flujos concretos entraron en él.
Con 358 flujos benignos genuinos en el día del SSH y 203 flujos de ataque en el
día web, un macro-F1 sin intervalo no permitiría distinguir una vista de otra.

La exactitud no se usa, por el motivo que ilustra la
sección~\ref{sec:reg-cifrado}. En el capítulo~\ref{ch:generalizacion} se añaden
el AUC y el error de calibración: cuando se compara entre dominios, una métrica
que dependa de un umbral fijo puede ocultar por completo lo que el modelo ha
aprendido.

\subsection*{Las tres vistas}

Se comparan tres representaciones del mismo flujo, siempre las mismas:

\begin{itemize}
  \item \emph{Payload}: histograma de los 256 valores de byte de la carga útil
        agregada de la conexión, más su entropía normalizada.
  \item Metadatos de flujo: cuatro escalares de la conexión, que son el número
        de paquetes, los bytes totales, la longitud de la secuencia y la
        duración, con $\log(1+x)$ para las colas pesadas.
  \item Conducta: los metadatos anteriores más tres características de ráfaga,
        el número de conexiones que el mismo origen ha abierto en el segundo,
        los cinco segundos y los treinta segundos previos, calculadas de forma
        causal, mirando solo al pasado, como haría un detector en línea.
\end{itemize}

Que sean siempre las mismas no es una obviedad. Una versión anterior de este
trabajo comparaba la vista conductual en un día contra los metadatos básicos en
otro, ambas rotuladas igual, y eso exageraba la diferencia entre regímenes.
Sobre ese error se vuelve en la sección~\ref{sec:reg-claro}.

\section{Los cuatro regímenes: dónde vive la firma}
\label{sec:regimenes-descriptivo}

Antes de medir, describir. Los seis días cubren cuatro maneras
cualitativamente distintas de que un ataque se haga visible en el tráfico. Esta
sección las caracteriza inspeccionando los datos, sin clasificador de por medio:
son hechos sobre los ataques, no sobre los modelos.

\subsection{Régimen cifrado: la firma vive en la repetición}
\label{sec:reg-cifrado}

El día \texttt{Wednesday-14-02-2018} contiene un ataque de fuerza bruta contra
SSH: 94.207 flujos de ataque desde \texttt{13.58.98.64} contra
\texttt{172.31.69.25:22}, entre las 14:01 y las 15:33 hora local, sobre un total
de 2.202.231 flujos.

El tráfico va cifrado y los bytes que transporta son, por diseño,
indistinguibles del ruido: la entropía de la carga útil es de aproximadamente
$7.4$ tanto en el ataque como en el tráfico legítimo. Lo que sí distingue al
ataque es que se repite, y la figura~\ref{fig:rafaga} muestra hasta qué punto:
la mediana de conexiones por origen en ventanas de cinco segundos es de 86 para
el atacante frente a 1 para el tráfico legítimo. Un ataque de fuerza bruta es,
por definición, un patrón de repetición.

\paragraph{El espejismo del 99.96\%.} Este día sirve además para mostrar por qué
la exactitud no aparece como métrica en el resto del capítulo. Un primer modelo
entrenado sobre el histograma de bytes y la entropía alcanzaba una exactitud del
$99.95\%$ en validación cruzada y del $99.96\%$ sobre el conjunto de prueba.
Auditado contra el tráfico SSH benigno, que es el único que un detector
desplegado encontraría a diario, su tasa de falsos positivos era del 96\%. El
modelo no había aprendido a reconocer el ataque: había aprendido a reconocer el
protocolo SSH, que ese día aparecía casi siempre asociado al ataque. Las dos
cifras son correctas y son inútiles.

La figura~\ref{fig:espejismo} pone las dos juntas, que es la única manera de
verlo. Por separado, la primera parece un resultado excelente y la segunda un
modelo distinto.

\begin{figure}[htbp]
  \centering
  \includegraphics[width=\textwidth]{F12_espejismo}
  \caption[El espejismo del 99.96\,\%]{El mismo modelo, medido de dos maneras.
  A la izquierda, la exactitud sobre un conjunto benigno compuesto en su
  práctica totalidad por otros protocolos. A la derecha, los falsos positivos
  sobre tráfico SSH benigno real, que es el que un detector desplegado
  encontraría.}
  \label{fig:espejismo}
\end{figure}

\subsection{Régimen en claro: la firma vive en los bytes}
\label{sec:reg-claro}

El día \texttt{Thursday-22-02-2018} contiene ataques web contra una aplicación
DVWA: 203 flujos repartidos en fuerza bruta (142), XSS (42) e inyección SQL
(19), desde \texttt{18.218.115.60} contra \texttt{172.31.69.28:80}.

Aquí el contenido malicioso está literalmente en los bytes: las cadenas
\texttt{union select}, \texttt{<script>} o la repetición de \texttt{POST
login.php} son el ataque, no un rastro de él. Y la conducta, que en el SSH era
la señal más nítida, aquí no existe: el contraste de ráfaga es de 1 conexión del
ataque frente a 3 del benigno, es decir, el ataque abre menos conexiones que el
tráfico normal. Es un ataque de bajo volumen.

\subsubsection*{El tipo de ataque vive en la secuencia, no en la composición}

El día web permite además una pregunta que los otros no admiten, porque es el
único con varios tipos de ataque distintos sobre el mismo servicio. En lugar de
separar ataque de benigno, se puede pedir al modelo que distinga entre los tres
tipos: fuerza bruta (142 flujos), XSS (42) e inyección SQL (19).

El resultado, en la tabla~\ref{tab:multitipo}, es el único de todo el trabajo en
que la representación del contenido importa por sí misma. La secuencia de bytes
identifica el tipo con un macro-F1 de $0.9505$, mientras el histograma de los
mismos bytes se queda en $0.7462$ y los metadatos de flujo en $0.7695$.

\begin{table}[htbp]
  \centering
  \caption[Identificación del tipo de ataque web]{Macro-F1 \emph{out-of-fold} al
  distinguir entre los tres tipos de ataque del día web, y \emph{recall} de la
  clase minoritaria, que es la inyección SQL con 19 flujos.}
  \label{tab:multitipo}
  \begin{tabular}{lrr}
    \toprule
    \textbf{Vista} & \textbf{Macro-F1} & \textbf{Recall SQLi} \\
    \midrule
    Payload, histograma & 0.7462 & 0.2632 \\
    Payload, secuencia  & 0.9505 & 0.8421 \\
    Metadatos de flujo  & 0.7695 & 0.2632 \\
    \bottomrule
  \end{tabular}
\end{table}

La diferencia se concentra en la clase minoritaria. El histograma y los
metadatos recuperan el 26.3\% de las inyecciones SQL, es decir, cinco de
diecinueve; la secuencia recupera el 84.2\%, dieciséis de diecinueve. El motivo
es que la inyección SQL y el ataque de fuerza bruta usan aproximadamente los
mismos caracteres, y su composición de bytes es casi idéntica y solo
el orden los distingue. La cadena \texttt{union select} no es un conjunto de
letras, es una secuencia.

Es un resultado modesto por el tamaño de la muestra, con diecinueve ejemplos en
la clase minoritaria, pero vale la pena retenerlo porque matiza todo el capítulo:
cuando la comparación entre vistas satura, la comparación entre
\emph{representaciones de la misma vista} todavía puede discriminar.

\subsection{Régimen volumétrico: la firma vive en el agregado}
\label{sec:reg-volumetrico}

Tres días cubren este régimen, y juntos muestran algo que ninguno enseñaría por
separado: el contraste de ráfaga se desploma tres órdenes de magnitud a medida
que el ataque se distribuye.

\begin{table}[htbp]
  \centering
  \caption[Contraste de ráfaga por día]{Mediana de conexiones abiertas por un
  mismo origen en ventanas de cinco segundos, para el ataque y para el tráfico
  legítimo del mismo día.}
  % @fuente [x for d in ["16-02 DoS-Hulk (1 origen)", "15-02 GoldenEye+Slowloris (1+1)", "20-02 DDoS-LOIC (10 origenes)", "14-02 SSH-Bruteforce", "02-03 Bot (beaconing)", "22-02 Web (BF/XSS/SQLi)"] for x in R.RAFAGA_MEDIANA[d]]
  \label{tab:rafaga-contraste}
  \begin{tabular}{lrrr}
    \toprule
    \textbf{Día y ataque} & \textbf{Ataque} & \textbf{Benigno} & \textbf{Contraste} \\
    \midrule
    16-02 DoS-Hulk (1 origen)            & 11.626 & 3 & $\times$3.875 \\
    15-02 GoldenEye + Slowloris (1+1)    &    104 & 3 & $\times$35 \\
    20-02 DDoS-LOIC (10 orígenes)        &     39 & 3 & $\times$13 \\
    14-02 SSH-Bruteforce                 &     86 & 1 & $\times$86 \\
    02-03 Bot (\emph{beaconing})         &     10 & 2 & $\times$5 \\
    22-02 Web (BF / XSS / SQLi)          &      1 & 3 & $\times$0 \\
    \bottomrule
  \end{tabular}
\end{table}

La lectura de la tabla~\ref{tab:rafaga-contraste} es que la ráfaga no es una
propiedad de los ataques volumétricos sino de los ataques concentrados. Un DDoS
repartido entre diez orígenes diluye por construcción la señal que un DoS de un
solo origen exhibe de forma masiva. El día web cierra el rango por el otro
extremo, con contraste nulo.

\begin{figure}[htbp]
  \centering
  \includegraphics[width=\textwidth]{F17_rafaga_distribucion}
  \caption[Distribución de la ráfaga por origen]{Conexiones abiertas por un mismo
  origen en ventanas de cinco segundos, en el día del SSH-Bruteforce. La mediana
  del atacante es 86; la del tráfico legítimo, 1.}
  \label{fig:rafaga}
\end{figure}

\subsection{Régimen de periodicidad: la firma vive en el intervalo}
\label{sec:reg-periodicidad}

El día \texttt{Friday-02-03-2018} contiene la actividad de una botnet (Ares):
142.921 flujos. Su firma no es el volumen, ya que su contraste de ráfaga es de
apenas $\times$5, sino la regularidad temporal. La figura~\ref{fig:beacon} lo
mide. El \emph{beacon} late con dos intervalos discretos, en torno a $0.5$ y
$2.0$ segundos, que concentran el 91.1\% de sus intervalos entre conexiones en
dos bandas de 50 milisegundos. El tráfico benigno del mismo día es una curva
suave decreciente: en la banda de $0.50$ a $0.55$ segundos cae el 41.8\% del
ataque frente al 1.3\% del benigno, un contraste de $\times$32.

Es el régimen que motivó una vista adicional, porque ninguna de las tres vistas
del protocolo calcula el intervalo entre conexiones. Sobre ello vuelve la
sección~\ref{sec:sintesis-regimenes}.

\section{La medición uniforme: todo satura}
\label{sec:saturacion}

Descritos los regímenes, la pregunta natural es cuál de las tres vistas gana en
cada uno. La tabla~\ref{tab:saturacion} la responde, y la respuesta es que la
pregunta no se puede contestar con estos datos.

\begin{table}[htbp]
  \centering
  \caption[Macro-F1 intra-día en los seis días]{Macro-F1 \emph{out-of-fold} de
  las mismas tres vistas, con el mismo protocolo y la misma métrica, en los seis
  días. Media de diez submuestreos.}
  % @fuente [R.SATURACION_INTRADIA[d][v][0] for d in DIAS for v in VISTAS]
  \label{tab:saturacion}
  \begin{tabular}{llll}
    \toprule
    \textbf{Día y régimen} & \textbf{Payload} & \textbf{Metadatos} & \textbf{Conducta} \\
    \midrule
    14-02 cifrado       & 0.9842 & 0.9906 & 0.9985 \\
    22-02 en claro      & 0.9734 & 0.9186 & 0.9823 \\
    16-02 volumétrico   & 1.0000 & 0.9715 & 0.9999 \\
    15-02 volumétrico   & 0.9998 & 0.9621 & 0.9967 \\
    20-02 vol.\ diluido & 1.0000 & 0.9987 & 0.9984 \\
    02-03 periodicidad  & 0.9990 & 0.9981 & 0.9983 \\
    \bottomrule
  \end{tabular}
\end{table}

El \emph{payload} no baja de $0.9734$ en ningún día y es el mejor o empata en
cuatro de los seis. La conducta no baja de $0.9823$. La única vista que falla
de forma apreciable son los metadatos de flujo básicos, y solo en tres días. La
figura~\ref{fig:saturacion} lo muestra como lo que es: una rejilla uniformemente
alta, sin contraste que leer.

\begin{figure}[htbp]
  \centering
  \includegraphics[width=\textwidth]{F16_matriz_regimen_vista}
  \caption[La evaluación intra-día está saturada]{La misma medición de la
  tabla~\ref{tab:saturacion}, con lo único que distingue unas celdas de otras
  escrito debajo de cada una: cuánto sobrevive de esa vista cuando el atacante
  camufla los primeros bytes del flujo.}
  \label{fig:saturacion}
\end{figure}

Vale la pena decir con claridad lo que esto derriba, porque era el marco con el que
se empezó este trabajo. La hipótesis inicial era que cada régimen escondería su
firma en una vista distinta, con lo que ninguna ganaría en todos y el enfoque
híbrido quedaría justificado por esa vía. Medido con una sola métrica y un solo
protocolo, ese marco no se sostiene: en el régimen volumétrico, donde se
esperaba que ganase el agregado, gana el \emph{payload}; y en el día de la
botnet, donde se había afirmado que ninguna vista mide la firma, las tres pasan
de $0.998$.

\subsection*{La única diferencia real}

Lo único que se mueve de forma estadísticamente sólida al cambiar de régimen es
el desplome de los metadatos per-flujo. La figura~\ref{fig:inversion} lo aísla:
entre el día cifrado y el día en claro, el \emph{payload} y la conducta apenas
se mueven y sus intervalos se solapan, mientras los metadatos caen 7.2 puntos y
son la única diferencia cuyos intervalos no se solapan.

\begin{figure}[htbp]
  \centering
  \includegraphics[width=\textwidth]{F19_inversion_web_ssh}
  \caption[Lo único que se invierte son los metadatos per-flujo]{Las mismas tres
  vistas en los dos regímenes opuestos, con barras de error al 95\%.}
  \label{fig:inversion}
\end{figure}

Merece una nota de honestidad metodológica, porque una versión anterior de esta
figura exageraba el efecto. Comparaba la vista conductual en el lado del SSH
contra los metadatos básicos en el lado web, ambas rotuladas «metadatos»,
simplemente porque el día web nunca se había evaluado con la vista conductual.
La caída aparente era de $1.000$ a $0.929$, y buena parte de esa caída no
era el cambio de régimen sino el cambio de características. Medido como es
debido, el efecto sigue existiendo y sigue siendo el único significativo, pero
es la mitad de grande.

\section{Por qué satura: una hipótesis y su refutación}
\label{sec:por-que-satura}

Un resultado tan uniforme pide una explicación. La primera que se propuso en
este trabajo resultó ser falsa, y merece contarse junto con el test que la
tumbó, porque el mecanismo correcto solo aparece al descartarla.

\subsection*{La hipótesis: el atajo es la identidad del host}

Cada uno de estos días contiene un único host atacante, con la excepción del
DDoS del 20 de febrero, que tiene diez. Ese host genera decenas o centenares de
miles de flujos que comparten todo lo que un host comparte consigo mismo: la
misma pila TCP, el mismo cliente, las mismas cabeceras y el mismo ritmo.

Bajo esas condiciones cabe sospechar que la tarea que se le plantea al
clasificador no es distinguir un ataque del tráfico normal sino, en la práctica,
distinguir las conexiones de ese host de las de los demás. Esa tarea es fácil
desde cualquier representación, lo que explicaría la saturación. Sería el
fenómeno que Engelen y sus coautores describen para CIC-IDS2017
\cite{Engelen2021} bajo el nombre de \emph{shortcut learning}: el modelo aprende
un atajo correlacionado con la etiqueta en lugar del concepto que la etiqueta
pretende nombrar.

La sospecha era razonable, y afectaba especialmente a la vista conductual,
porque las características de ráfaga cuentan conexiones por origen y con un
único atacante podrían estar funcionando como un identificador de ese host más
que como una descripción de su conducta.

\subsection*{El test: apartar un atacante entero}

La hipótesis es contrastable, pero no en un día con un solo atacante. El DDoS
del 20 de febrero sí lo permite, porque tiene diez hosts atacantes con unos
29.000 flujos cada uno. Si el éxito se debiera a memorizar el host, evaluar
sobre un atacante nunca visto debería hundirlo.

El protocolo aparta en cada pliegue un host atacante completo y, además, un
grupo disjunto de hosts benignos, para que tampoco la clase negativa pueda
memorizarse por host. Se entrena con 5.000 flujos por clase y se evalúa sobre
2.000, y se compara contra un control de exactamente los mismos tamaños con
partición aleatoria por flujo: sin él, una caída podría deberse al tamaño del
entrenamiento y no al cambio de host.

Se añade un segundo control, este positivo: una cuarta vista deliberadamente
tramposa formada por los cuatro octetos de la dirección IP de origen, que no
puede sino memorizar el host. Si el protocolo es sensible al efecto que pretende
medir, esa vista debe acertar en la partición aleatoria y desplomarse al apartar
un atacante. Sin esa comprobación, un resultado negativo no significaría nada.

El algoritmo~\ref{alg:loao} lo detalla, incluidos los dos controles.

\begin{algorithm}[htbp]
\caption{Leave-one-attacker-out con control negativo y positivo}
\label{alg:loao}
\begin{algorithmic}[1]
\Require Hosts atacantes $H_A$; hosts benignos $H_B$; vistas $V$ más la vista
         tramposa de octetos de la IP
\Ensure Macro-F1 por vista, apartando un atacante, y su control aleatorio
\State Repartir $H_B$ en $|H_A|$ grupos disjuntos al azar
\For{$i = 1$ \textbf{hasta} $|H_A|$}
  \State $\mathit{test} \gets$ 2000 flujos del atacante $h_i$ y 2000 del grupo benigno $i$
  \State $\mathit{train} \gets$ 5000 flujos de $H_A \setminus \{h_i\}$ y 5000 de los demás grupos
  \ForAll{vista $v \in V$}
    \State entrenar en $\mathit{train}$, predecir $\mathit{test}$, anotar macro-F1
  \EndFor
\EndFor
\Statex
\Comment{Control negativo: mismos tamaños, partición aleatoria por flujo}
\For{$i = 1$ \textbf{hasta} $|H_A|$}
  \State $\mathit{train}, \mathit{test} \gets$ muestra aleatoria de igual tamaño, ignorando el host
  \ForAll{vista $v \in V$}
    \State entrenar, predecir y anotar macro-F1
  \EndFor
\EndFor
\Statex
\Comment{Control positivo: la vista de octetos debe hundirse en el bucle 1}
\end{algorithmic}
\end{algorithm}

\begin{table}[htbp]
  \centering
  \caption[Leave-one-attacker-out sobre el DDoS del 20-02]{Macro-F1 al apartar un
  host atacante completo, frente a un control aleatorio de idéntico tamaño.
  Media y desviación sobre los diez pliegues. La última fila no es una vista
  candidata: es un identificador de host puro, incluido para verificar que el
  test detecta la memorización cuando existe.}
  % @fuente [x for v in ["payload (histograma + entropia)", "metadatos de flujo", "conducta (metadatos + rafaga)", "[control] octetos de la IP de origen"] for x in R.LOAO_20_02[v]]
  \label{tab:loao}
  \begin{tabular}{lccc}
    \toprule
    \textbf{Vista} & \textbf{Control} & \textbf{Atacante nuevo} & \textbf{Dif.} \\
    \midrule
    Payload (histograma + entropía) & 1.0000 $\pm$ 0.0000 & 1.0000 $\pm$ 0.0000 & $+0.0000$ \\
    Metadatos de flujo              & 0.9991 $\pm$ 0.0005 & 0.9992 $\pm$ 0.0006 & $+0.0001$ \\
    Conducta (metadatos + ráfaga)   & 0.9986 $\pm$ 0.0006 & 0.9990 $\pm$ 0.0007 & $+0.0004$ \\
    \midrule
    \emph{[control] octetos de la IP} & 0.9999 $\pm$ 0.0001 & 0.9332 $\pm$ 0.2000 & $-0.0667$ \\
    \bottomrule
  \end{tabular}
\end{table}

\subsection*{El resultado: la hipótesis era falsa}

La tabla~\ref{tab:loao} no deja margen. Apartar un host atacante entero no
degrada ninguna de las tres vistas: las diferencias son de $+0.0000$,
$+0.0001$ y $+0.0004$, es decir, nulas. Lo que las vistas aprenden no es la
identidad del host.

El control positivo merece leerse con cuidado, porque su media esconde el dato
interesante. Cae de $0.9999$ a $0.9332$, pero con una desviación de
$0.2000$: el efecto no está repartido sino concentrado en un solo
pliegue. El motivo es que nueve de los diez atacantes tienen direcciones en el
rango \texttt{18.2xx.x.x} y el tráfico benigno está en \texttt{172.31.x.x}, así
que apartar un atacante deja otros nueve del mismo rango y el primer octeto
sigue separando. Es la fuga de direcciones IP del
capítulo~\ref{ch:generalizacion}, medida aquí a nivel de subred y no de host.

El pliegue decisivo es el de \texttt{52.14.136.135}, la única atacante fuera de
ese rango. Ahí el control positivo se hunde a $0.3333$, es decir, predice una
sola clase, mientras las tres vistas reales siguen dando $1.0000$, $0.9987$
y $0.9990$. En el único pliegue donde se demuestra que un atajo por identidad
de host falla, las vistas reales aguantan intactas. El test es sensible, y aun
así no detecta memorización de host en ninguna de ellas.

\subsection*{Lo que sí explica la saturación}

Descartado el host, queda la herramienta. Los diez atacantes de ese día ejecutan
el mismo programa, y todos los ataques de todos los días de este conjunto están
generados por una única herramienta cada uno. Lo que las vistas aprenden es esa
herramienta y las condiciones de la captura, no el ataque como categoría.

La confirmación viene del mismo día por otra vía. Camuflar los primeros bytes
del flujo hunde el \emph{payload} de $1.0000$ a $0.0108$
(sección~\ref{sec:motivo-equivocado}). Los dos experimentos se complementan de
forma precisa: el leave-one-attacker-out dice que no es el host, y la evasión
dice que es la herramienta.

El alcance de este test es limitado. Este test se puede hacer en un solo día
de los seis, porque es el único con más de un atacante, y en los otros cinco la
hipótesis del host sigue siendo, estrictamente, no contrastada. Y como los diez
atacantes comparten herramienta, el test distingue memorizar el host de aprender
la herramienta o la conducta, pero no puede distinguir aprender la herramienta
de aprender el ataque. Para eso hacen falta la evasión y la transferencia entre
dominios.

\paragraph{Una reserva que el test levanta.} La sospecha sobre la vista
conductual, que la ráfaga fuese un identificador de host disfrazado, queda
desmentida con datos: medida contra un atacante que nunca vio, la conducta rinde
igual que contra atacantes conocidos, con $0.9990$ frente a $0.9986$. Cuenta
conexiones por origen, sí, pero lo que captura es un patrón de comportamiento
que se repite en hosts distintos.

\paragraph{Qué se puede afirmar y qué no.} De la tabla~\ref{tab:saturacion} se
sigue que las tres vistas contienen información suficiente para separar estos
ataques de este tráfico benigno. No se sigue que el problema esté resuelto, ni
que una vista sea mejor que otra, ni que un detector entrenado así funcionaría
en una red real. Un macro-F1 de $0.99$ en este montaje no mide capacidad de
detección.

\paragraph{Un resultado anterior que queda explicado.} Durante el trabajo se
contrastaron varias arquitecturas del estado del arte, un \emph{transformer}
sobre bytes y una red de atención cruzada, y ninguna mejoró los resultados,
algunas a un coste computacional entre dos y cincuenta veces mayor. En su
momento se interpretó como un límite de los datos. La
tabla~\ref{tab:saturacion} lo explica mejor: no mejoraban porque no quedaba nada
que mejorar. El cuello de botella nunca estuvo en la capacidad del modelo.

\section{Aciertan, pero por el motivo equivocado}
\label{sec:motivo-equivocado}

Si las tres vistas aciertan, la pregunta útil deja de ser cuánto aciertan y pasa
a ser por qué. Ahí sí se separan, porque no todas aciertan por la misma razón ni
con la misma solidez.

\subsection*{El payload cifrado lee el banner, no el cifrado}

El $0.9842$ del \emph{payload} en el día del SSH no procede de leer contenido
cifrado, que sería imposible. Procede del \emph{banner} en claro que precede a
la negociación de la sesión y que identifica al cliente. El atacante emite dos
banners distintos, y el modal es literalmente
\texttt{SSH-2.0-paramiko\_2.0.0}, frente a los 81 de los clientes legítimos como
PuTTY, libssh2 u OpenSSH. El análisis de saliencia lo confirma: el 31.1\% de la
importancia atribuida se concentra en los primeros 48 bytes, la zona del banner,
y no en el cuerpo cifrado.

Es una firma real, pero de la herramienta y no del ataque. El mismo patrón se
repite en los días volumétricos, donde lo que el modelo aprende son las
cabeceras y los \emph{User-Agent} característicos de Hulk, GoldenEye y LOIC.

\subsection*{La prueba: camuflar los primeros bytes}

La forma de comprobarlo es directa. Si el modelo se apoya en los primeros bytes
del flujo, sobrescribirlos con los de un cliente legítimo debería hundirlo, y a
un atacante real le costaría muy poco hacerlo.

Eso es exactamente lo que ocurre. El experimento completo se desarrolla en la
sección~\ref{sec:evasion}, tabla~\ref{tab:evasion-cap5}, porque pertenece al
capítulo de las perturbaciones. Basten aquí las dos cifras de titular: el
\emph{recall} del \emph{payload} en el día de la botnet cae de $1.0000$ a
$0.0025$, y en el DDoS a $0.0108$. El \emph{payload} deja de funcionar
mientras el ataque sigue funcionando igual de bien, porque el atacante no ha
cambiado su técnica, solo su saludo.

\paragraph{La excepción confirma la regla, y hubo que medirla dos veces.} El día
web es el único que aguanta, y el camino hasta esa cifra merece contarse. El
primer experimento sobrescribía 160 bytes y el \emph{recall} se desplomaba a
$0.0164$, lo que parecía confirmar que también allí la firma era una huella de
herramienta. Al inspeccionar los flujos modificados quedó claro que el
experimento estaba mal planteado: en 160 bytes cabe la inyección entera, de modo
que sobrescribirlos no evade la detección, destruye el ataque. Repetido con un
prefijo de 24 bytes, que no llega a la carga maliciosa, el \emph{payload}
aguanta en $0.8852$. Su firma sí es contenido genuino.

\paragraph{La conducta es invariante, y no por mérito.} Las vistas de metadatos
y de conducta no cambian un solo valor bajo este ataque, sencillamente porque no
leen el contenido. Es una ventaja estructural, no una virtud del modelo, y hay
que presentarla como tal: dice que esta forma concreta de evasión no las afecta,
no que sean robustas frente a cualquier evasión. Un atacante que ralentice
deliberadamente su fuerza bruta ataca la ráfaga igual de barato.

\begin{figure}[htbp]
  \centering
  \includegraphics[width=\textwidth]{F37_interarribo_beacon}
  \caption[El beacon de la botnet, medido]{Distribución de los intervalos entre
  conexiones consecutivas del mismo origen. El ataque concentra el 91.1\% de sus
  esperas en dos bandas de 50 ms; el tráfico benigno del mismo día es una curva
  suave.}
  \label{fig:beacon}
\end{figure}

\section{El límite del enfoque: lo que el payload no puede ver}
\label{sec:limite}

Los apartados anteriores comparan vistas que aciertan. Este trata de los casos
en que una de ellas no tiene nada que mirar, que es una situación
cualitativamente distinta de acertar poco.

\subsection*{Un ataque sin payload en el propio día del SSH}

El día 14 de febrero contiene un segundo ataque, de fuerza bruta contra FTP, que
no aparece en ninguna tabla de este capítulo porque no genera un solo flujo con
carga útil analizable. El atacante (\texttt{18.221.219.4}) dirige 162.801
conexiones a \texttt{172.31.69.25:21}, de las cuales 134.945 quedan en estado
\texttt{REJ} con cero bytes, porque el servicio FTP de la víctima estaba caído
durante la captura y rechazó los intentos. Las restantes tampoco llegan a
exportar datos de aplicación. El neto es cero flujos FTP con \emph{payload}.

La causa es más concreta de lo que parece. Es una circunstancia de esta captura, un
servicio caído, y no una propiedad del protocolo FTP. Lo que generaliza no es el
mecanismo sino la consecuencia, y esa consecuencia reaparece por otras vías: el
SlowHTTPTest del 16 de febrero también deja cero flujos con \emph{payload}, en
su caso porque un ataque de agotamiento lento mantiene conexiones abiertas sin
llegar a intercambiar datos.

\subsection*{Infiltration: el 94.4\% de una intrusión sin un byte que analizar}

El caso más severo es el día de la Infiltration, y llega por la vía del
resultado negativo. Tras comprometer una máquina, el atacante ejecuta una fase
de reconocimiento desde dentro de la red. De las 193.097 conexiones salientes de
la víctima, 182.323, el 94.4\%, no llevan un solo byte de aplicación: 158.474
quedan en estado \texttt{S0}, que es un SYN sin respuesta, y 21.069 en
\texttt{REJ}. Zeek no identifica servicio en 183.907 de ellas. Son un barrido de
puertos contra 612 destinos.

La figura~\ref{fig:ceguera} reparte esas conexiones según lleven o no carga
útil. El contraste dentro del mismo día es lo que cierra el argumento: las 5
conexiones al canal de mando y control (\texttt{13.58.225.34:31337}) sí llevan
\emph{payload}, con 195.816 bytes de aplicación y el 100\% de ellas con
contenido, y son perfectamente visibles. El análisis de contenido ve la parte
pequeña de la intrusión y es ciego a la grande.

No es que el \emph{payload} funcione peor: es que no puede ver. Una fase entera
de una intrusión real es estructuralmente invisible para un enfoque basado en
contenido, y ninguna mejora del modelo lo arregla, porque no hay datos que
mejorar. La literatura recoge el mismo problema al reetiquetar CIC-IDS2017 bajo
la categoría «\emph{Attempted}» \cite{Engelen2021}
(sección~\ref{sec:sota-datasets}).

Este resultado es, con diferencia, el argumento más fuerte a favor de combinar
vistas, y es el único de todo el capítulo que no depende de una comparación de
métricas.

\begin{figure}[htbp]
  \centering
  \includegraphics[width=\textwidth]{F35_ceguera_infiltracion}
  \caption[La ceguera del payload ante la Infiltration]{Reparto de las 193.097
  conexiones salientes de la víctima según lleven o no carga útil de
  aplicación. El 94.4\% no lleva un solo byte que analizar.}
  \label{fig:ceguera}
\end{figure}

\section{Síntesis: la pregunta correcta}
\label{sec:sintesis-regimenes}

Este capítulo empezó con la intención de mostrar qué vista gana en cada régimen
y termina mostrando que dentro del día no gana ninguna en particular, porque el
montaje experimental no lo permite. Queda recoger qué sigue en pie.

Los cuatro regímenes son reales. La firma del SSH-Bruteforce vive en la
repetición, con 86 conexiones por ventana frente a 1; la de los ataques web en
los bytes; la del DoS en el agregado; y la de la botnet en el intervalo, con el
91.1\% de sus esperas en dos bandas de 50 ms. Eso se comprueba inspeccionando el
tráfico y no depende de ningún clasificador.

Lo que no se sostiene es la asignación de una vista ganadora a cada régimen. Con
un solo host atacante por día, cualquier vista dispone de un atajo suficiente, y
las diferencias entre ellas caen dentro del ruido.

Hay una excepción a la saturación, y apunta al capítulo siguiente. Ninguna de
las tres vistas calcula el intervalo entre conexiones, que es donde vive la
firma de la botnet. El día 02-03 se detecta igualmente al $0.999$, por la
huella HTTP del canal de mando, pero se detecta por el motivo equivocado, y por
eso es también el día en que la evasión resulta más devastadora, de $1.0000$ a
$0.0025$. La saturación y la fragilidad son dos caras de lo mismo.

La pregunta correcta es, por tanto, otra. Si dentro del día todas las vistas
aciertan, lo que hay que preguntar no es cuál acierta más, sino cuál sigue
acertando cuando cambian las condiciones: cuando el atacante camufla sus bytes,
cuando el modelo se aplica a otro día o a otro conjunto de datos, y cuando hay
que fijar un umbral de decisión sin ver antes los datos de destino. Ese es el
contenido del capítulo~\ref{ch:generalizacion}.
